% ============================================
% Johns Hopkins Letter Template
% Uses the built-in 'letter' class
% ============================================

\documentclass[11pt,letterpaper]{letter}

\usepackage{graphicx}
\usepackage{eso-pic}
\usepackage{geometry}
\usepackage{microtype}
\usepackage[utf8]{inputenc}
\usepackage[hidelinks]{hyperref}

% ----- Page Setup -----
\geometry{left=25mm,right=25mm,top=25mm,bottom=25mm}

% ----- Logo Placement (foreground, first page only) -----
\newcommand{\PutJHULogoFG}[1]{%
  \AddToShipoutPictureFG*{%
    \AtPageUpperLeft{%
      \hspace{10mm}%
      \raisebox{-38mm}{%
        \includegraphics[width=6cm]{#1}%
      }%
    }%
  }%
}

% ----- Sender Information -----
\signature{Decory Edwards}
\address{Department of Economics \\ Johns Hopkins University \\ Baltimore, MD 21218 \\ \texttt{dedwar65@jh.edu}}

% ----- Recipient Info -----
\begin{document}
\PutJHULogoFG{Images/jhulogo.png} % show logo only on first page

\begin{letter}{Hiring Committee \\
The Heritage Foundation \\
214 Massachusetts Ave NE \\
Washington, DC 20002}

\opening{Dear Hiring Committee,}

I am writing to express my interest in the Senior Policy Analyst -- Macroeconomics and Public Finance position at The Heritage Foundation. I am a Ph.D.\ candidate in the Department of Economics at Johns Hopkins University (advisors: Christopher D.\ Carroll, Nicholas W.\ Papageorge, and Stelios Fourakis), and I expect to receive my degree in May 2026. My primary specializations are macroeconomic modeling, public finance, and computational economics.

My research agenda focuses on the ability of macroeconomic models to generate empirically realistic distributions of wealth and income over the life cycle. A key driver of that work is modeling heterogeneity in the rate of return on assets, and understanding how return dispersion contributes to portfolio accumulation across households. These themes connect directly to Heritage’s focus on evaluating federal tax and spending programs through rigorous economic reasoning, transparent modeling, and evidence-based assessments of long-run fiscal sustainability.

In my job-market paper, I calibrate a life-cycle heterogeneous-agent model to match wealth moments from the Survey of Consumer Finances by allowing households to differ in their portfolio returns. A central component of the paper compares the effects of a capital-income tax and a wealth tax under return heterogeneity—quantifying implications for savings behavior, inequality, economic efficiency, and welfare. This analysis provides a disciplined, data-driven framework for evaluating how tax structures influence economic incentives and long-run macroeconomic outcomes, aligning with Heritage’s commitment to advancing policies that promote economic opportunity and sound fiscal governance.

In a related research project using the Health and Retirement Study, I examine how trust influences both income growth and asset returns—finding a hump-shaped relationship for returns and characterizing the distribution of the persistent component of returns to net worth. These experiences have equipped me to: (i) frame research strategies and design modeling approaches; (ii) execute econometric and simulation analysis with large micro-data sets; and (iii) translate quantitative findings into clear, rigorous written deliverables suitable for policy audiences.

I am particularly motivated by Heritage’s mission to evaluate economic policy through the principles of opportunity, fiscal responsibility, and long-run prosperity. I would welcome the opportunity to contribute my analytical rigor, modeling experience, and policy-oriented research to the Macroeconomics and Public Finance team.

Thank you for considering my application.

\closing{Sincerely,}

\end{letter}
\end{document}
